\documentclass[12pt,aspectratio=43]{beamer}

%%%%%%%%%%%%%%%%%%%%%%%%%%%%%%%%%%%%%%%%%%%%%%%%%%
%   Paquetes de configuración basica del texto   %
%%%%%%%%%%%%%%%%%%%%%%%%%%%%%%%%%%%%%%%%%%%%%%%%%%
\usepackage[spanish,es-nodecimaldot]{babel}
\usepackage{babelbib}
\usepackage[T1]{fontenc}
\usepackage{listings}
\usepackage{multirow}
\usepackage{hyperref}
\usepackage{hyperxmp}
\usepackage{ifthen}

%%%%%%%%%%%%%%%%%%%%%%%%%%%%%%%%%%%%%%%%%%%%%%%%%%
%                   [listings]                   %
%   Configuración de vizualizacion del código    %
%%%%%%%%%%%%%%%%%%%%%%%%%%%%%%%%%%%%%%%%%%%%%%%%%%
\lstset{
	backgroundcolor=\color{lightgray!10},
	basicstyle=\ttfamily,
	extendedchars=true,
	showspaces=false,
	showstringspaces=false,
	captionpos=b,
	keywordstyle=\bfseries\color{cyan},
	commentstyle=\color{gray},
	stringstyle=\color{orange},
	escapeinside={!>}{<!},
	language=[LaTeX]TeX }

%%%%%%%%%%%%%%%%%%%%%%%%%%%%%%%%%%%%%%%%%%%%%%%%%%
%     Paquetes de configuración graficos del     %
%                   documento                    %
%%%%%%%%%%%%%%%%%%%%%%%%%%%%%%%%%%%%%%%%%%%%%%%%%%
\usepackage{graphicx}
\usepackage{tikz}
\usetikzlibrary{babel,shadows}
\usepackage{xcolor}
\usepackage{pdfpages}
\usepackage{animate}

\usepackage{ifxetex}
\ifxetex
	\usepackage[no-math]{fontspec}
	\setmainfont{AncizarSans}[
		Path = ../Fuentes/,
		Extension = .otf,
		BoldFont = *-B,
		ItalicFont = *-I,
		BoldItalicFont = *-BI ]
	\setmonofont{UbuntuMono}[
		Path = ../Fuentes/,
		Extension = .ttf,
		BoldFont = *-B,
		ItalicFont = *-I,
		BoldItalicFont = *-BI ]
	\newcommand{\lmr}{\fontfamily{lmr}\selectfont}
	\newcommand{\lmss}{\fontfamily{lmss}\selectfont}
	\newcommand{\lmtt}{\fontfamily{lmtt}\selectfont}
\fi

%%%%%%%%%%%%%%%%%%%%%%%%%%%%%%%%%%%%%%%%%%%%%%%%%%
%           Configuraciones de Beamer            %
%%%%%%%%%%%%%%%%%%%%%%%%%%%%%%%%%%%%%%%%%%%%%%%%%%
\definecolor{Igreen}{RGB}{148,180,59}
\definecolor{Ired}{RGB}{166,24,49}

\setbeamertemplate{navigation symbols}{}

\setbeamertemplate{itemize items}{\color{Igreen}\raisebox{.45ex}{\rule{.6ex}{.6ex}}}

\usefonttheme{professionalfonts}
\usefonttheme{serif}
\setbeamercolor{frametitle}{fg=Ired,bg=Igreen!50}
\setbeamercolor{alerted text}{fg=Ired}
\setbeamertemplate{enumerate item}{\color{Igreen}\bfseries\arabic{enumi}.}
\setbeamertemplate{enumerate subitem}{\color{Igreen}\bfseries\arabic{enumi}.\arabic{enumii}.}

\makeatletter
\newcommand{\ifratio}[2]{
	\ifthenelse
	{\lengthtest{\beamer@paperwidth=16cm} \AND \lengthtest{\beamer@paperheight=9cm}}
	{#1}
	{#2} }
\makeatother

\setbeamertemplate{title page}{
	\begin{tikzpicture}[overlay,remember picture]
	\path (current page.north west) ++(0.5cm,-0.5cm) node[below right] {\includegraphics[height=1.25cm]{../Escudo_UN}};
	\path (current page.north east) ++(-0.5cm,-0.5cm) node[below left] {\parbox[c][1.25cm][c]{\widthof{\Large \insertsubtitle}}{\Large \insertsubtitle}};
	\end{tikzpicture}
	
	\vfill
	\begin{center}
		{\Huge\inserttitle}\\
		\bigskip
		\insertauthor
	\end{center}
	\vfill
}

%%%%%%%%%%%%%%%%%%%%%%%%%%%%%%%%%%%%%%%%%%%%%%%%%%
%         Información de la presentación         %
%%%%%%%%%%%%%%%%%%%%%%%%%%%%%%%%%%%%%%%%%%%%%%%%%%

\title{Elementos Básicos de Edición}
\subtitle{Curso Libre de {\lmr\LaTeX}}
\author{Joar Esteban Buitrago Carrillo}
\institute{Universidad Nacional de Colombia}
\date{}
\logo{\includegraphics[height=0.75cm]{../Escudo_UN}}
\hypersetup{
	pdfcopyright={Esta obra está bajo una licencia de Creative Commons Reconocimiento-CompartirIgual 4.0 Internacional.},
	pdflicenseurl={https://creativecommons.org/licenses/by-sa/4.0/deed.es}
}

\begin{document}
\begin{frame}[plain]
\titlepage
\end{frame}

{
\setbeamercolor{background canvas}{bg=}
\begin{frame}
\ifratio
	{\includepdf[pages=2]{../Licencia_Diapositiva.pdf}}
	{\includepdf[pages=1]{../Licencia_Diapositiva.pdf}}
\end{frame}
}

\section{Modos de Procesamiento de {\LaTeX}}

\newcommand{\blindtext}{Y, viéndole don Quijote de aquella manera, con muestras de tanta tristeza, le dijo: Sábete, Sancho, que no es un hombre más que otro si no hace más que otro. Todas estas borrascas que nos suceden son señales de que presto ha de serenar el tiempo y han de sucedernos bien las cosas; porque no es posible que el mal ni el bien sean durables, y de aquí se sigue que, habiendo durado mucho el mal, el bien está ya cerca.}

\begin{frame}{Modos de Procesamiento de {\lmr\LaTeX}}{}
{\lmr\LaTeX} posee diferentes modos para el procesamiento del texto.\\[1em]

\begin{itemize}[<+->]
	\item \alert<1>{\bf Normal}
	\item \alert<2>{\bf Izquierda-Derecha} I-D
	\item \alert<3>{\bf Matemático}
\end{itemize}
\end{frame}

\subsection{Modo Normal}
\begin{frame}{Modo Normal}{Modos de Procesamiento de {\lmr\LaTeX}}
En el modo normal el texto se maneja de forma natural, separando los párrafos al llegar al margen derecho.\pause\\[1em]

Por ejemplo:\\[1em]

\begin{minipage}[c]{\linewidth}
	\lmr\blindtext
\end{minipage}
\end{frame}

\subsection{Modo Izquierda-Derecha}
\begin{frame}{Modo Izquierda-Derecha}{Modos de Procesamiento de {\lmr\LaTeX}}
En el modo Izquierda-Derecha nunca se separan los párrafos al llegar al margen derecho, lo cual genera un error en la visualización.\pause\\[1em]

Por ejemplo:\\[1em]

\begin{minipage}[c]{\linewidth}\hfuzz=2000pt % Suprime mostrar el error
	\mbox{\lmr\blindtext}
\end{minipage}\pause\\[1em]

{\lmr\LaTeX} indicara que hace falta espacio para acomodar el texto. Así:\\[1em]

\texttt{Overfull \textbackslash{}hbox (1848.40244pt too wide) in paragraph at lines 166--166}
\end{frame}

\subsection{Modo Matemático}
\begin{frame}{Modo Matemático}{Modos de Procesamiento de {\lmr\LaTeX}}
El modo matemático se genera al entrar en un entorno o comando matemático. Funciona de manera similar que el modo I-D, con excepción del cambio de fuente y la exclusión de todos los espacios en blanco.\pause\\[1em]

Por ejemplo:\\[-0.5em]

\begin{minipage}[c]{\linewidth}\hfuzz=2000pt % Lo mismo aqui
	\[\blindtext\]
\end{minipage}\pause\\[1em]

En este caso {\lmr\LaTeX} también indicara que hace falta espacio para acomodar el texto. Así:\\[1em]

\texttt{Overfull \textbackslash{}hbox (1590.5191pt too wide) detected at line 184}
\end{frame}

\section{Listas}
\begin{frame}{Listas}{}
Uno de los elementos de edición más importantes de cualquier trabajo escrito son las listas.\pause\\[1em]

Como cualquier otro editor de texto {\lmr\LaTeX} también posee varias formas de incluir listas en el documento.\pause\\[1em]

En {\lmr\LaTeX} existen 3 tipos básicos de listas, las listas \alert{numeradas}, las \alert{no numeradas} y las \alert{descripciones}.
\end{frame}

\subsection{Listas Numeradas}
\begin{frame}[fragile]{Listas Numeradas}{Listas}
Para escribir una lista \alert{\bf enumerada} utilizamos el entorno {\tt enumerate}, y para cada uno de los elementos de la lista utilizamos el comando {\tt item}.\pause\\[1em]

Así:
\begin{columns}
	\begin{column}{0.5\linewidth}
		{\small
\begin{lstlisting}
\begin{enumerate}
  \item Elemento 1.
  \item Elemento 2.
  \item ...
\end{enumerate}
\end{lstlisting} }
		{\color{Ired} Código}
	\end{column}
	\vline
	\begin{column}{0.5\linewidth}
		{\lmr{\small
				% Por el comportamiento de Beamer no puedo poner el entorno normal
				\hspace{2ex}
				\parbox{3cm}{
					1. Elemento 1.\\
					2. Elemento 2.\\
					3. ... } }}\\
		\hfill{\color{Ired} Vista}
	\end{column}
\end{columns}
\end{frame}

\subsection{Listas no Numeradas}
\begin{frame}[fragile]{Listas no Numeradas}{Listas}
Para escribir una lista \alert{\bf no enumerada}, o por viñetas, utilizamos el entorno {\tt itemize}, y para cada uno de los elementos de la lista utilizamos el comando {\tt item}.\pause\\[1em]

Así:
\begin{columns}
	\begin{column}{0.5\linewidth}
		{\small
\begin{lstlisting}
\begin{itemize}
  \item Elemento 1.
  \item Elemento 2.
  \item ...
\end{itemize}
\end{lstlisting} }
		{\color{Ired} Código}
	\end{column}
	\vline
	\begin{column}{0.5\linewidth}
		{\lmr{\small
				% Por el comportamiento de Beamer no puedo poner el entorno normal
				\hspace{2ex}
				\parbox{3cm}{
					\raisebox{.45ex}{\rule{.6ex}{.6ex}} Elemento 1.\\
					\raisebox{.45ex}{\rule{.6ex}{.6ex}} Elemento 2.\\
					\raisebox{.45ex}{\rule{.6ex}{.6ex}} ... } }}\\
		\hfill{\color{Ired} Vista}
	\end{column}
\end{columns}
\end{frame}

\subsection{Descripciones}
\begin{frame}[fragile]{Descripciones}{Listas}
Para escribir una \alert{\bf descripción} utilizamos el entorno {\tt description}, y para cada uno de los elementos de la lista utilizamos el comando {\tt item} con un argumento opcional.\pause\\[1em]

Así:
\begin{columns}
	\begin{column}{0.5\linewidth}
		{\small
\begin{lstlisting}
\begin{itemize}
  \item[Desc 1:] Texto.
  \item[Desc 2:] Texto.
  \item[Desc 3:] ...
\end{itemize}
\end{lstlisting} }
		{\color{Ired} Código}
	\end{column}
	\vline
	\begin{column}{0.5\linewidth}
		{\lmr{\small
				% Por el comportamiento de Beamer no puedo poner el entorno normal
				\hspace{2ex}
				\parbox{4cm}{
					\textbf{Desc 1:} Texto.\\
					\textbf{Desc 2:} Texto.\\
					\textbf{Desc 3:} ... } }}\\
		\hfill{\color{Ired} Vista}
	\end{column}
\end{columns}
\end{frame}

\subsection{Anidación de Listas}
\begin{frame}[fragile]{Anidación de Listas}{Listas}
Para anidar listas es muy simple, ya que solamente se requiere de abrir y cerrar entornos.\pause\\[1em]

Así:
\begin{columns}
	\begin{column}{0.5\linewidth}
		{\small
\begin{lstlisting}
\begin{enumerate}
  \item Texto.
  \begin{enumerate}
    \item Anidado.
  \end{enumerate}
\end{enumerate}
\end{lstlisting} }
		{\color{Ired} Código}
	\end{column}
	\vline
	\begin{column}{0.5\linewidth}
		{\lmr{\small
				% Por el comportamiento de Beamer no puedo poner el entorno normal
				\hspace{2ex}
				\parbox{4cm}{
					1. Texto.\\
					\hspace*{2em}a) Anidado.} }}\\
		\hfill{\color{Ired} Vista}
	\end{column}
\end{columns}
\end{frame}

\section{Cajas}

\subsection{mbox y fbox}

\subsubsection{mbox}

\subsubsection{fbox}

\subsection{makebox y framebox}

\subsubsection{makebox}

\subsubsection{framebox}

\subsection{parbox}

\section{mini-páginas}

\section{Notas al pie}

\section{Notas al Margen}

\section{El entorno {\tt verbatim} y su paquete}
\end{document}